\section{Preliminaries}
\label{sec:preliminaries}

We reuse several basic operations and their verified properties from
companion articles:

\begin{itemize}
  \item \textbf{Modular Arithmetic} \cite{mata2026modulo}: Division,
    modulo, quotient invariance, mod idempotence.
  \item \textbf{Lists} \cite{mata2026lists}: Size, append, sum, slicing,
    tail shift.
  \item \textbf{Prime Utilities:} Primorial computation, primality
    testing, and finite divisor search, specified below.
\end{itemize}

\subsection{Key Definitions}
\label{sec:key-definitions}

Let $L = [p_1, p_2, \dots, p_k] \in \mathbb{N}^k$ be a non-empty list of
primes, with $p_i > 1$ for all $i$.

We define the \textbf{primorial} of a list of primes as the product of all
primes in the list:

\begin{equation*}
\operatorname{primorial}(L) = \prod_{i=1}^{k} p_i
\end{equation*}

A number $n$ is \textbf{prime} exactly when it is greater than $1$ and no
integer in $[2,n)$ divides it:

\begin{equation*}
\operatorname{isPrime}(n) \;:\Longleftrightarrow\; n>1\ \land\
\forall d\in[2,n),\ \operatorname{mod}(n,d)\ne0.
\end{equation*}

\subsection{Finite Prime Operations and Their Specifications}
\label{sec:finite-prime-operations}

The article needs only finite search and the definitions below; it does
not assume an external enumeration of all primes. Primorial is structural:

\begin{equation*}
\begin{aligned}
\operatorname{primorial}([]) &:= 1, \\
\operatorname{primorial}(p::P) &:= p\cdot\operatorname{primorial}(P).
\end{aligned}
\end{equation*}

For $n>1$ and $2\leq s\leq n$, \texttt{findSmallestDivisor(n,s)} tests
$s,s+1,\ldots,n$ in order and returns the first divisor. It terminates
because $n$ itself divides $n$. Its specification is therefore

\begin{equation*}
\begin{aligned}
d &= \operatorname{findSmallestDivisor}(n,s) \\
&\implies s\leq d\leq n\ \land\ \operatorname{mod}(n,d)=0 \\
&\phantom{\implies}\land\ \forall e\in[s,d),\ \operatorname{mod}(n,e)\ne0.
\end{aligned}
\end{equation*}

The recursive scan proves this by induction on $n-s$: either $s$ divides
$n$, or the same claim is inherited from the call beginning at $s+1$.
Thus, when $s=2$, $d=n$ is equivalent to the prime definition in
Section~2.1; if $d<n$, $d$ is the least non-trivial divisor. The
divisor-range postcondition is encoded in
\href{https://github.com/thiagomata/prime-numbers/blob/master/src/main/scala/v1/chapter5/prime/Prime.scala}{\texttt{Prime::findSmallestDivisor}}.
Its minimality and its equivalence with an empty divisor range are
machine-checked by \texttt{Prime::assertFindSmallestDivisorMinimality} and
\texttt{Prime::assertFindSmallestDivisorEquivNoDivisorInRange} in the same
source.

For a list $P$ of positive integers, \texttt{isCoprime(n,P)} means that no
member of $P$ divides $n$; it is a direct recursion over the list. This is
the only list-coprimality meaning used below.
